El análisis de algoritmos permite estudiar cómo crece el costo de una solución
cuando aumenta el tamaño de la entrada. En este informe se estudian dos
problemas fundamentales: el ordenamiento de arreglos unidimensionales y la
multiplicación de matrices cuadradas. Se comparan implementaciones de
MergeSort, QuickSort, PatienceSort y \texttt{std::sort}, además de los métodos
Naive y Strassen, relacionando sus tiempos de ejecución y uso de memoria con
las complejidades teóricas descritas en la literatura
\cite{taocv3,algorithms_erickson}. El objetivo es identificar qué tan bien las
mediciones experimentales reflejan el comportamiento asintótico esperado y
cómo influyen en ellas el tamaño, la estructura y la distribución de los datos.
